%% toc: P15 161 161.1
% Convergence in pth mean from convergence in probability, with source forms of the measure foundations and of the conditional zero-one law (Durrett, author version 5, pp. 1--4, 10, 13, 15, 16, 24, 117, 244--245, 247--250, 456).
\chapter{Mean-Power Convergence from Convergence in Probability}
\label{sec:p15-w02c-p15-014-probability-mean}

When a sequence of random variables converges in probability, the
question whether it also converges in \(p\)th mean is a question about
uniform integrability. This chapter states that link exactly, for every
exponent \(0<p<\infty\), with the integrability hypotheses that each
direction needs, and shows by models that none of those hypotheses can
be removed. It then records source forms in the measure foundations
(the definitions of a sigma-field and of a field, subadditivity in
covering form, the names for the Borel sets, the generality of measurable maps and Borel functions of
several random variables, the half-open rectangles
and Dynkin's pi-lambda theorem, and point masses as laws), the source forms of uniform integrability and
of its moment criterion, and source discussions of conditional limits
that complete the presentations of
\cref{sec:p15-measurable-probability,sec:p15-uniform-integrability-and-conditional-limits}.

\emph{Purpose.} To give a correct characterization of \(p\)th-mean
convergence among sequences that converge in probability, valid
whether or not the variables have finite \(p\)th moments.

\emph{Inputs.} A probability space \((\Omega,\mathcal F,\mathbb P)\);
finite real random variables \(X_n\) (\(n\geq1\)) and \(X\) on it with
\(X_n\to X\) in probability; an exponent \(0<p<\infty\). The modes are
those of \cref{def:p15-w02b-p15-013-modes} and uniform integrability is
that of \cref{def:p15-uniform-integrability}.

\emph{Conclusions.} \(X_n\to X\) in \(L^p\) exactly when some tail of
\((|X_n-X|^p)\) is uniformly integrable
(\cref{thm:p15-w02c-p15-014-probability-mean}~(1)). If every \(X_n\)
lies in \(L^p\), this is equivalent to uniform integrability of
\((|X_n|^p)_{n\geq1}\) (item~(2)); if \(X\) lies in \(L^p\), it is
equivalent to uniform integrability of some tail of \((|X_n|^p)\)
(item~(3)). Without either hypothesis, uniform integrability of
\((|X_n|^p)_{n\geq1}\) together with \(X_n\in L^p\) for every \(n\) is
sufficient but not necessary (item~(4)).

\section{The characterization}
\label{sec:p15-w02c-p15-014-characterization}

\begin{hbconvention}{def:p15-w02c-p15-014-reading}{Reading the limit of \(p\)th means}
For \(0<p<\infty\), \cref{def:p15-w02b-p15-013-modes}~(4) says that
\(X_n\to X\) in \(L^p\) when \(\mathbb E|X_n-X|^p\to0\). The numbers
\(a_n=\mathbb E|X_n-X|^p\) lie in \([0,\infty]\), and \(a_n\to0\)
means: for every \(\varepsilon>0\) there is \(N\) with
\(a_n\leq\varepsilon\) for all \(n\geq N\). Finitely many of the
\(a_n\) may therefore be infinite, and neither \(X_n\) nor \(X\) is
required to lie in \(L^p\). Uniform integrability is defined in
\cref{def:p15-uniform-integrability} only for families of integrable
variables; a family containing a variable of infinite mean is not
uniformly integrable. Durrett's definition, which applies to any
collection, agrees with it and excludes such families as well
(\cref{thm:p15-w02c-p15-014-ui-arbitrary}). For \(N\geq1\), the \emph{tail} of a sequence
\((V_n)_{n\geq1}\) from \(N\) is the family \((V_n)_{n\geq N}\).
Throughout the chapter
\[
 V_n\defeq|X_n-X|^p,\qquad W_n\defeq|X_n|^p,\qquad W\defeq|X|^p .
\]
These are nonnegative, finite and measurable.
\end{hbconvention}

\begin{hbtheorem}{thm:p15-w02c-p15-014-probability-mean}{proved=annex:p15-w02c-p15-014-prob-mean}{From convergence in probability to convergence in \(p\)th mean}
Let \(0<p<\infty\) and let \(X_n\) (\(n\geq1\)) and \(X\) be finite real
random variables on one probability space with \(X_n\to X\) in
probability.
\begin{enumerate}
\item \(X_n\to X\) in \(L^p\) if and only if there is \(N\geq1\) such
      that \(\mathbb E|X_n-X|^p<\infty\) for every \(n\geq N\) and the
      tail \((|X_n-X|^p)_{n\geq N}\) is uniformly integrable. If
      \(\mathbb E|X_n-X|^p<\infty\) for every \(n\), the condition is
      that \((|X_n-X|^p)_{n\geq1}\) is uniformly integrable. The
      implication from \(L^p\) convergence to uniform integrability of
      a tail holds without the hypothesis of convergence in probability.
\item Suppose \(X_n\in L^p\) for every \(n\). Then \(X_n\to X\) in
      \(L^p\) if and only if \((|X_n|^p)_{n\geq1}\) is uniformly
      integrable, and in that case \(X\in L^p\).
\item Suppose \(X\in L^p\). Then \(X_n\to X\) in \(L^p\) if and only if
      there is \(N\geq1\) such that \(X_n\in L^p\) for every \(n\geq N\)
      and \((|X_n|^p)_{n\geq N}\) is uniformly integrable.
\item The condition ``\(X_n\in L^p\) for every \(n\) and
      \((|X_n|^p)_{n\geq1}\) is uniformly integrable'' implies
      \(X_n\to X\) in \(L^p\), but \(L^p\) convergence does not imply it.
      Let \(Y_p\) be a finite positive random variable with
      \(\mathbb E Y_p^{\,p}=\infty\); one exists on the uniform space of
      \cref{thm:p15-lebesgue-unit-interval} for every \(p\).
      \begin{enumerate}
      \item If \(X_n=X=Y_p\) for every \(n\), then \(X_n\to X\) in
            probability and in \(L^p\), while no \(X_n\) lies in
            \(L^p\) and \((|X_n|^p)\) is not uniformly integrable.
      \item If \(X=0\), \(X_1=Y_p\) and \(X_n=0\) for \(n\geq2\), then
            \(X\in L^p\) and \(X_n\to X\) in probability and in
            \(L^p\), while \(X_1\notin L^p\); the tail from \(N=2\)
            satisfies item~(3).
      \end{enumerate}
\end{enumerate}
\end{hbtheorem}

Item~(2) is the corrected form of the probability-to-mean comparison:
the equivalence between \(L^p\) convergence and uniform integrability of
\((|X_n|^p)_{n\geq1}\) holds under the standing hypothesis that every
\(X_n\) lies in \(L^p\), as in \cref{thm:p15-w02b-p15-013-vitali}, and
item~(4) shows that the hypothesis cannot be moved into the conclusion.
Item~(1) is the form that needs no moment hypothesis at all, and
item~(3) the form in which only the limit is assumed to have a finite
\(p\)th moment. At \(p=1\), item~(2) is Durrett's Theorem~4.6.3, which
assumes \(\mathbb E|X_n|<\infty\) for all \(n\) and then proves
uniform integrability, \(L^1\) convergence and convergence of
\(\mathbb E|X_n|\) to \(\mathbb E|X|<\infty\) equivalent
\cite[Theorem~4.6.3, p.~245]{DurrettPTEAuthor2019}; in this handbook it
is \cref{thm:p15-vitali-convergence}. Items~(1), (3) and~(4) and the
range \(p\neq1\) are not in the source and are proved in the annex.

\begin{hbmethod}{def:p15-w02c-p15-014-method}{Deciding \(p\)th-mean convergence after convergence in probability}
\emph{Input:} finite real random variables \(X_n\), a candidate limit
\(X\) on the same space, and \(0<p<\infty\).
\begin{enumerate}
\item Establish \(X_n\to X\) in probability, for example from
      almost-sure convergence (\cref{thm:p15-w02b-p15-013-comparison}~(1)).
      If \((X_n)\) converges in probability to some \(X'\), any
      \(L^p\) limit equals \(X'\) almost surely
      (\cref{thm:p15-w02b-p15-013-comparison}~(2) and
      \cref{thm:p15-w02b-p15-013-as-probability}~(2)), so only \(X'\)
      needs to be tested.
\item Determine which moment hypothesis holds. If
      \(\mathbb E|X_n|^p<\infty\) for every \(n\), use item~(2); if only
      \(\mathbb E|X|^p<\infty\) is known, use item~(3); otherwise use
      item~(1) with \(V_n=|X_n-X|^p\).
\item Test uniform integrability of the family named by that item:
      by domination \(|X_n|\leq Z\) with \(Z\in L^p\); by a bound
      \(\sup_n\mathbb E|X_n|^r<\infty\) for some \(r>p\)
      (\cref{thm:p15-w02b-p15-013-vitali}); or by computing
      \(\sup_{n\geq N}\mathbb E(U_n;U_n>M)\) directly for the family
      \((U_n)\) in question.
\item Conclude \(L^p\) convergence if the family is uniformly
      integrable. If it is not, and in item~(1) no tail is uniformly
      integrable, conclude that \(X_n\) does not converge in \(L^p\)
      to any random variable.
\end{enumerate}
\emph{Side conditions.} All variables are finite, real and defined on
one space of total mass one; for an infinite measure the step from
convergence everywhere to convergence in measure already fails
(\cref{thm:p15-w02b-p15-013-finite-measure}). A bound on
\(\mathbb E|X_n|^p\) alone is not a test of uniform integrability
(\cref{thm:p15-w02b-p15-013-countermodels}~(1)).
\end{hbmethod}

\begin{hbexample}{def:p15-w02c-p15-014-worked}{Spikes of two heights}
On the uniform space of \cref{thm:p15-lebesgue-unit-interval}, with
sample space \((0,1)\) and probability \(\mathbb U\), let
\(c_n>0\) and \(X_n=c_n\mathbf1_{(0,1/n)}\), with \(X=0\). For
\(\varepsilon>0\), \(\mathbb U(|X_n|>\varepsilon)\leq1/n\to0\), so
\(X_n\to0\) in probability, and every \(X_n\) lies in \(L^p\) because
\(\mathbb E|X_n|^p=c_n^p/n\). The variable \(|X_n|^p\) is
\(c_n^p\mathbf1_{(0,1/n)}\), so for \(M>0\)
\[
 \mathbb E\bigl(|X_n|^p;|X_n|^p>M\bigr)=
 \begin{cases} c_n^p/n,& c_n^p>M,\\ 0,& c_n^p\leq M.\end{cases}
\]
\emph{Height \(c_n=n^{1/(2p)}\).} Then \(c_n^p=\sqrt n\) and the
display equals \(n^{-1/2}\) when \(n>M^2\) and zero otherwise, so the
supremum over \(n\) is at most \(1/M\), which tends to zero. The family
\((|X_n|^p)\) is uniformly integrable, and item~(2) gives
\(X_n\to0\) in \(L^p\); directly,
\(\mathbb E|X_n|^p=n^{-1/2}\to0\).

\emph{Height \(c_n=n^{1/p}\).} Then \(c_n^p=n\) and the display equals
one for every \(n>M\). For every \(N\) and \(M\) the supremum over
\(n\geq N\) is one, so no tail of \((|X_n|^p)=(|X_n-0|^p)\) is
uniformly integrable, and item~(1) shows that \(X_n\) does not converge
to zero in \(L^p\); indeed \(\mathbb E|X_n|^p=1\) for every \(n\). This
is model~(1) of \cref{thm:p15-w02b-p15-013-countermodels}.
\end{hbexample}

\begin{hbwarning}{def:p15-w02c-p15-014-limits}{Hypotheses that cannot be dropped}
\begin{enumerate}
\item The hypothesis \(X_n\in L^p\) for every \(n\) is needed for the
      necessity half of item~(2): in both models of item~(4),
      \(X_n\to X\) in \(L^p\) and \((|X_n|^p)_{n\geq1}\) is not
      uniformly integrable.
\item Convergence in probability is needed for the sufficiency half of
      item~(1). Let \(\zeta\) be the fair sign of
      \cref{thm:p15-w02b-p15-013-countermodels}~(4), \(X_n=-\zeta\) and
      \(X=\zeta\). Then \(|X_n-X|^p=2^p\) for every \(n\), a constant
      family and hence uniformly integrable, while
      \(\mathbb E|X_n-X|^p=2^p\) does not tend to zero.
\item Uniform integrability is not implied by a bound on the \(p\)th
      moments: the second height in
      \cref{def:p15-w02c-p15-014-worked} has
      \(\mathbb E|X_n|^p=1\) for every \(n\).
\end{enumerate}
\end{hbwarning}

\section{Source forms in the measure foundations}
\label{sec:p15-w02c-p15-014-foundations}

Durrett's opening definitions and his appendix on the extension theorem
contain variants of the definitions of a sigma-algebra and of an
algebra, the covering form of subadditivity, notation for the Borel sets,
a stated reason for the generality of measurable maps with Durrett's
theorem on Borel functions of several random variables, point masses as
laws, and an example of a pi-system with the attribution of the
pi-lambda theorem, which \cref{sec:p15-measurable-probability} replaces by neighbouring
forms. They are given here in the source's form, with the bridge to the
handbook's form proved.

\begin{hblemma}{thm:p15-w02c-p15-014-sigma-field-variant}{proved=annex:p15-w02c-p15-014-field}{Durrett's definitions of a sigma-field and of a field}
\par\nobreak
Let \(\Omega\) be a nonempty set and \(\mathcal F,\mathcal A\subseteq2^\Omega\).
\begin{enumerate}
\item \(\mathcal F\) is a sigma-algebra in the sense of
      \cref{def:p15-sigma-algebra} if and only if \(\mathcal F\) is
      nonempty and satisfies
      \begin{enumerate}
      \item if \(A\in\mathcal F\), then \(A^c\in\mathcal F\);
      \item if \(A_i\in\mathcal F\) for each member of a countable
            sequence, then \(\bigcup_iA_i\in\mathcal F\).
      \end{enumerate}
      Countable means finite or countably infinite.
\item \(\mathcal A\) is an algebra in the sense of
      \cref{def:p15-sigma-algebra} if and only if \(\mathcal A\) is
      nonempty and \(A,B\in\mathcal A\) implies \(A^c\in\mathcal A\)
      and \(A\cup B\in\mathcal A\). Every algebra is then closed under
      intersections of two members.
\item The empty collection satisfies the closure conditions of~(1)
      and of~(2) and is neither a sigma-algebra nor an algebra.
\end{enumerate}
\end{hblemma}

The right-hand side of~(1) is Durrett's definition: he asks for a
nonempty collection with the two closure properties and does not list
\(\Omega\in\mathcal F\) \cite[sec.~1.1, p.~1]{DurrettPTEAuthor2019}.
Closure under countable intersections follows by the identity
\(\bigcap_iA_i=(\bigcup_iA_i^c)^c\), and Durrett gives his reason for not
including that property in the definition: the shorter list is easier to
check \cite[p.~1]{DurrettPTEAuthor2019}. To verify that a collection
is a sigma-field one then checks two closure properties and exhibits one
member. In the handbook the derived closure properties are
\cref{thm:p15-generated-events}. Durrett then names the pair
\((\Omega,\mathcal F)\), taken without the probability, a
\emph{measurable space}, and glosses the name: it is a space on which a
measure can be put \cite[sec.~1.1, p.~1]{DurrettPTEAuthor2019}. The
gloss matches the order of his definitions, since the measure he defines
next is a function on \(\mathcal F\) and needs nothing beyond the pair;
\cref{def:p15-sigma-algebra} uses the term in the same sense.

The right-hand side of~(2), without the word nonempty, is the
definition Durrett gives on his way from a semialgebra to a
sigma-algebra; he calls such a collection an \emph{algebra}, or
\emph{field}, and obtains closure under intersection from
\(A\cap B=(A^c\cup B^c)^c\) \cite[sec.~1.1, p.~4]{DurrettPTEAuthor2019}.
Item~(3) shows that nonemptiness must be read into his condition for the
two definitions to agree. The word \emph{field} here names a collection of sets
closed under the finite set operations, and a sigma-field is a field
closed also under countable unions. It has nothing to do with the
algebraic structure of the same name. Durrett notes that every
sigma-algebra is an algebra and gives the finite and cofinite subsets of
the integers as an algebra that is not a sigma-algebra
\cite[Example~1.1.6, p.~4]{DurrettPTEAuthor2019}; the same example on
the positive integers is \cref{def:p15-finite-cofinite-events}.
\index{field of sets}\index{algebra of sets}\index{sigma-field}

\begin{hblemma}{thm:p15-w02c-p15-014-covering-subadditivity}{proved=annex:p15-w02c-p15-014-covering}{Subadditivity in covering form}
\par\nobreak
Let \(\mu\) be a measure on \((\Omega,\mathcal F)\), and let \(A\) and
\(A_1,A_2,\ldots\) belong to \(\mathcal F\).
\begin{enumerate}
\item \emph{Subadditivity.} If \(A\subset\bigcup_{m=1}^\infty A_m\), then
      \[
       \mu(A)\leq\sum_{m=1}^\infty\mu(A_m),
      \]
      where the sum of the numbers \(\mu(A_m)\in[0,\infty]\) may be
      infinite.
\item The statement in~(1), quantified over all such \(A\) and
      \((A_m)\), is equivalent to the union form
      \(\mu(\bigcup_mA_m)\leq\sum_m\mu(A_m)\) of
      \cref{thm:p15-measure-laws}~(2), given monotonicity
      (\cref{thm:p15-measure-laws}~(1)).
\item The same bound holds for a finite cover
      \(A\subset A_1\cup\cdots\cup A_k\), with the sum over
      \(m\leq k\).
\end{enumerate}
\end{hblemma}

Durrett states the bound in the covering form of~(1) and names it
\emph{subadditivity}; it is item~(ii) of his Theorem~1.1.1, among the
consequences of the definition of a measure listed there with
monotonicity and the two continuity properties
\cite[Theorem~1.1.1, p.~2]{DurrettPTEAuthor2019}. He assumes throughout
that theorem that every set mentioned lies in \(\mathcal F\), and the
lemma keeps that hypothesis: for a set \(A\) outside \(\mathcal F\) the
number \(\mu(A)\) is not defined. His proof makes the cover disjoint
inside \(A\) and then uses countable additivity and monotonicity; the
annex proof follows that route, so it does not depend on the union form.
\Cref{thm:p15-measure-laws}~(2) is the case \(A=\bigcup_mA_m\), and the
covering form is the one used when a set is bounded by a cover that is
larger than the set, as in the following example.
\index{subadditivity!of a measure}

\begin{hbexample}{def:p15-w02c-p15-014-covering-example}{Covering the rationals by short intervals}
Let \(\lambda\) be Lebesgue measure on \((\mathbb R,\mathcal B(\mathbb R))\)
(\cref{thm:p15-lebesgue-unit-interval}), and let
\(q_1,q_2,\ldots\) list the rational numbers. The set \(\mathbb Q\) is a
countable union of closed singletons, so it is a Borel set. Fix
\(\varepsilon>0\) and put
\(A_m=(q_m-\varepsilon2^{-m-1},\,q_m+\varepsilon2^{-m-1}]\), an interval
of length \(\varepsilon2^{-m}\). Every rational is some \(q_m\) and lies
in \(A_m\), so \(\mathbb Q\subset\bigcup_mA_m\), and
\cref{thm:p15-w02c-p15-014-covering-subadditivity}~(1) gives
\[
 \lambda(\mathbb Q)\leq\sum_{m=1}^\infty\varepsilon2^{-m}=\varepsilon .
\]
As \(\varepsilon\) is arbitrary, \(\lambda(\mathbb Q)=0\). The cover is
larger than the set: \(\bigcup_mA_m\) contains the interval \(A_1\) of
positive length, so \(\lambda(\bigcup_mA_m)>0=\lambda(\mathbb Q)\), and the
union form applied to \(\bigcup_mA_m\) alone would bound only the
cover. The bound is strict for every such cover, since
\(\lambda(\mathbb Q)=0<\varepsilon\). If the intervals were taken of fixed
length \(\varepsilon\), the sum would be infinite and the bound would
carry no information; subadditivity bounds a set only by a cover whose
measures are summable.
\end{hbexample}

\begin{hbconvention}{def:p15-w02c-p15-014-borel-notation}{Durrett's names for the Borel sets}
Durrett writes \(\mathbf R^d\) for the set of vectors
\((x_1,\ldots,x_d)\) of real numbers and \(\mathcal R^d\) for its
Borel sets, the smallest sigma-field containing the open sets; when
\(d=1\) he drops the superscript, so that \((\mathbf R,\mathcal R)\) is
the real line with its Borel sets \cite[sec.~1.1, p.~3]{DurrettPTEAuthor2019}.
The smallest sigma-field containing a collection is the generated
sigma-algebra of \cref{thm:p15-generated-events}, and
\cref{def:p15-borel-sets} defines \(\mathcal B(\mathbb R^d)\) as the
sigma-algebra generated by the open sets of \(\mathbb R^d\). Hence
\[
 \mathbf R^d=\mathbb R^d,\qquad
 \mathcal R^d=\mathcal B(\mathbb R^d),\qquad
 \mathcal R=\mathcal B(\mathbb R).
\]
Each of Durrett's uses translates by this dictionary. His definition of
a random variable asks that \(X^{-1}(B)\in\mathcal F\) for every
\(B\in\mathcal R\) \cite[sec.~1.2, p.~10]{DurrettPTEAuthor2019}, which is
\cref{def:p15-measurable-map} with
\((S,\mathcal S)=(\mathbb R,\mathcal B(\mathbb R))\). A random vector is a
measurable map into \((\mathbf R^d,\mathcal R^d)\) with \(d>1\)
\cite[sec.~1.3, p.~15]{DurrettPTEAuthor2019}. In his Theorem~1.3.5 the
function \(f\) is measurable from \((\mathbf R^n,\mathcal R^n)\) to
\((\mathbf R,\mathcal R)\), and the proof of his Theorem~1.3.6 uses that
the open set \(\{x:x_1+\cdots+x_n<a\}\) lies in \(\mathcal R^n\)
\cite[p.~16]{DurrettPTEAuthor2019}. This handbook writes
\(\mathcal B(\mathbb R^d)\) throughout; the letter \(\mathcal R\) has
other meanings in other chapters of this Part. The starred form for the
extended line is \cref{def:p15-w02c-p15-002-extended-line}.
\index{Borel set!notation}
\end{hbconvention}

\begin{hbconvention}{def:p15-w02c-p15-014-target-generality}{Why measurable maps take values in an arbitrary measurable space}
\Cref{def:p15-measurable-map} defines measurability for a map into an
arbitrary measurable space \((S,\mathcal S)\), and only then specializes
to real random variables and random vectors. This follows Durrett, who
gives the reason for working at that level: most of what is to be shown
holds for random elements of any measurable space, with the same proofs,
which are sometimes easier there \cite[sec.~1.3, p.~15]{DurrettPTEAuthor2019}.
The test by a generating class
(\cref{thm:p15-measurability-criterion}) and the measurability of
composites (\cref{thm:p15-measurable-operations}) are examples: their
proofs use only inverse images and the sigma-algebra operations, and no
order, metric or coordinate structure of \(S\). Structure of
\(\mathbb R\) or \(\mathbb R^d\) enters only through the choice of a
generating class, as in \cref{thm:p15-borel-generators}.
\end{hbconvention}

\begin{hbcorollary}{thm:p15-w02c-p15-014-borel-functions}{proved=annex:p15-w02c-p15-014-borel-functions}{Borel functions of random variables}
\par\nobreak
Let \(n\geq1\) and let \(X_1,\ldots,X_n\) be real random variables on
\((\Omega,\mathcal F,\mathbb P)\).
\begin{enumerate}
\item The map \(V=(X_1,\ldots,X_n):\Omega\to\mathbb R^n\) is a random
      vector: for Borel sets \(A_1,\ldots,A_n\subseteq\mathbb R\),
      \[
       \{V\in A_1\times\cdots\times A_n\}=\bigcap_{i=1}^n\{X_i\in A_i\}
       \in\mathcal F,
      \]
      and the products \(A_1\times\cdots\times A_n\) of Borel sets
      generate \(\mathcal B(\mathbb R^n)\).
\item If \(f:(\mathbb R^n,\mathcal B(\mathbb R^n))\to
      (\mathbb R,\mathcal B(\mathbb R))\) is measurable, then
      \(f(X_1,\ldots,X_n)\) is a real random variable.
\item The converse of~(2) fails: \(f(X_1,\ldots,X_n)\) can be a random
      variable for a continuous \(f\) when no \(X_i\) is one.
\end{enumerate}
\end{hbcorollary}

Item~(2) is Durrett's Theorem~1.3.5, in his notation
\(f:(\mathbf R^n,\mathcal R^n)\to(\mathbf R,\mathcal R)\)
(\cref{def:p15-w02c-p15-014-borel-notation}), and item~(1) is the step of
his proof \cite[Theorem~1.3.5, p.~16]{DurrettPTEAuthor2019}. He
introduces the theorem as the result that shows why his composition
theorem, Theorem~1.3.4, was proved for measurable maps between arbitrary
measurable spaces \cite[p.~16]{DurrettPTEAuthor2019}. The reason is
visible in the proof: the inner map of \(f(X_1,\ldots,X_n)\) is the
random vector \(V\), valued in \(\mathbb R^n\), and the outer map \(f\)
is defined on \(\mathbb R^n\), so a composition theorem stated only for
real random variables and functions on \(\mathbb R\) would not apply.
In the handbook the composition theorem is the first sentence of
\cref{thm:p15-measurable-operations}, stated for arbitrary measurable
spaces as \cref{def:p15-w02c-p15-014-target-generality} explains. That
theorem asserts the conclusion of~(2) for continuous \(f\); item~(2)
extends it to every Borel function, which includes functions with
discontinuities, as in the next example. Sums of random variables, his
Theorem~1.3.6, are the case \(f(x)=x_1+\cdots+x_n\)
\cite[Theorem~1.3.6, p.~16]{DurrettPTEAuthor2019}.
\index{random variable!Borel function of}

\begin{hbexample}{def:p15-w02c-p15-014-ratio-example}{The ratio of two random variables}
Define \(f:\mathbb R^2\to\mathbb R\) by \(f(x_1,x_2)=x_1/x_2\) when
\(x_2\neq0\) and \(f(x_1,0)=0\). The function is discontinuous at every
point \((x_1,0)\) with \(x_1\neq0\), since \(|f(x_1,t)|=|x_1|/|t|\) is
unbounded as \(t\to0\), so the continuous case of
\cref{thm:p15-measurable-operations} does not apply. It is Borel: for
real \(a\),
\[
 \{f\leq a\}
 =\{x_2>0,\ x_1-ax_2\leq0\}\cup\{x_2<0,\ x_1-ax_2\geq0\}
  \cup\{x_2=0,\ 0\leq a\},
\]
because multiplying \(x_1/x_2\leq a\) by \(x_2\) keeps the inequality
when \(x_2>0\) and reverses it when \(x_2<0\). Each of the sets
\(\{x_2>0\}\), \(\{x_2<0\}\), \(\{x_2=0\}\), \(\{x_1-ax_2\leq0\}\) and
\(\{x_1-ax_2\geq0\}\) is open or closed, as a set where a continuous
function is positive, negative, zero, at most zero or at least zero,
and the last set in the display is \(\{x_2=0\}\) or empty. So
\(\{f\leq a\}\) is a Borel set, and since the rays \((-\infty,a]\)
generate \(\mathcal B(\mathbb R)\) (\cref{thm:p15-borel-generators}),
\(f\) is measurable by \cref{thm:p15-measurability-criterion}. For real
random variables \(X_1,X_2\),
\cref{thm:p15-w02c-p15-014-borel-functions}~(2) shows that the ratio
\(f(X_1,X_2)\), equal to \(X_1/X_2\) where \(X_2\neq0\) and to 0
elsewhere, is a random variable.
\end{hbexample}

The hypothesis of item~(2) is a sufficient condition only. Item~(3)
shows that measurability of \(f(X_1,\ldots,X_n)\) gives no information
about the \(X_i\), and the annex example has \(n=2\),
\(f(x_1,x_2)=x_1-x_2\) and \(X_1=X_2\) an indicator of a set that is not
an event.

\begin{hblemma}{thm:p15-w02c-p15-014-half-open-rectangles}{proved=annex:p15-w02c-p15-014-rectangles}{Half-open rectangles form a generating pi-system}
\par\nobreak
Let \(d\geq1\). For real vectors \(a=(a_1,\ldots,a_d)\) and
\(b=(b_1,\ldots,b_d)\) put
\[
 R(a,b)\defeq(a_1,b_1]\times\cdots\times(a_d,b_d],
\]
which is empty when \(b_j\leq a_j\) for some \(j\), and let
\(\mathcal H_d\) be the collection of all such sets.
\begin{enumerate}
\item \(\mathcal H_d\) is a pi-system: with coordinatewise maxima and
      minima \(a\vee a'\) and \(b\wedge b'\),
      \[
       R(a,b)\cap R(a',b')=R(a\vee a',\,b\wedge b').
      \]
\item \(\sigma(\mathcal H_d)=\mathcal B(\mathbb R^d)\).
\item Two probability measures on \((\mathbb R^d,\mathcal B(\mathbb R^d))\)
      that agree on every half-open rectangle agree on every Borel set.
\end{enumerate}
\end{hblemma}

Durrett gives the half-open rectangles as his example of a pi-system
when he introduces the pi-lambda theorem
\cite[app.~A.1, p.~456]{DurrettPTEAuthor2019}. He attributes that
theorem to Dynkin, calling it \emph{Dynkin's \(\pi\)--\(\lambda\)
theorem}, and presents it as the key to the uniqueness part of his
extension theorem and as a result used many times later in his book
\cite[app.~A.1, p.~456]{DurrettPTEAuthor2019}. His Theorem~A.1.4 states
that if \(\mathcal P\) is a pi-system and \(\mathcal L\) a lambda-system
containing \(\mathcal P\), then \(\sigma(\mathcal P)\subset\mathcal L\)
\cite[Theorem~A.1.4, p.~456]{DurrettPTEAuthor2019}; with his lambda-system
conditions (\(\Omega\in\mathcal L\), closure under proper differences
\(B-A\) for \(A\subset B\), closure under increasing unions), which are
those of \cref{def:p15-pi-lambda}, this is the first sentence of
\cref{thm:p15-measure-uniqueness}. Durrett gives no reference to a work
of Dynkin at that point, and the attribution is recorded here as his.
In this handbook a lambda-system is also called a \emph{Dynkin system}.
\index{pi-lambda theorem}\index{Dynkin's pi-lambda theorem}\index{Dynkin system}
The empty set must be counted among the half-open rectangles: if only
rectangles with \(a_j<b_j\) for every \(j\)
were allowed, the collection would not be closed under intersection,
since \((0,1]\cap(1,2]=\varnothing\) in dimension one. Item~(3) is the
use to which the example is put, through
\cref{thm:p15-measure-uniqueness}.

\begin{hbdefinition}{def:p15-w02c-p15-014-point-mass-law}{Point mass as a law, and discrete laws}
For \(c\in\mathbb R\), Durrett's \emph{point mass at} \(c\) is a
probability measure on \((\mathbb R,\mathcal B(\mathbb R))\): the
measure \(\delta_c\) that gives a Borel set measure one if it contains
\(c\) and zero otherwise, the Dirac probability measure of
\cref{def:p15-probability-density}. The \emph{distribution function}
of a probability measure \(\mu\) on the line is
\(x\mapsto\mu((-\infty,x])\); for the law of \(X\) it is the function
\(F_X\) of \cref{def:p15-cumulative-distribution}. His Example~1.2.8, titled
\emph{Point mass at~0}, specifies it through its distribution function
\[
 F(x)=1\quad\text{for }x\geq0,\qquad F(x)=0\quad\text{for }x<0
\]
\cite[Example~1.2.8, p.~13]{DurrettPTEAuthor2019}. He calls a
probability measure \(P\) on the line, or its distribution function,
\emph{discrete} if there is a countable set \(S\) with \(P(S^c)=0\), and
presents the point mass at~0 as the simplest discrete distribution
\cite[p.~13]{DurrettPTEAuthor2019}. The term keeps this meaning later:
in his discussion of weak convergence the laws \(\mu_n\) are point masses
at \(1/n\) \cite[sec.~3.2, p.~117]{DurrettPTEAuthor2019}, the example
treated in \cref{thm:p15-weak-convergence-boundaries}~(1).

In \cref{def:p15-pushforward-law} the phrase \emph{point mass at}
\(x\) names a number, \(p_X(x)=\mathbb P(X=x)\), the value of the law
of \(X\) at the singleton \(\{x\}\). The two uses are tied by
\cref{thm:p15-w02c-p15-014-point-mass-law}: the law of \(X\) is the
point mass \(\delta_c\) in Durrett's sense exactly when the number
\(p_X(c)\) is one, and a discrete law is the weighted sum of the
measures \(\delta_x\) with the numbers \(p_X(x)\) as weights. In this
handbook \emph{point mass} without qualification keeps the numerical
meaning of \cref{def:p15-pushforward-law}; the measure is written
\(\delta_c\).
\index{point mass!as a law}\index{distribution!discrete}
\end{hbdefinition}

\begin{hblemma}{thm:p15-w02c-p15-014-point-mass-law}{proved=annex:p15-w02c-p15-014-point-mass}{Point-mass laws and discrete laws}
\par\nobreak
Let \(c\in\mathbb R\), and let \(X\) be a real random variable on
\((\Omega,\mathcal F,\mathbb P)\) with law \(\mathbb P_X\) and point
masses \(p_X(x)=\mathbb P(X=x)\).
\begin{enumerate}
\item \(\delta_c\) is a probability measure on
      \((\mathbb R,\mathcal B(\mathbb R))\) with distribution function
      \(F_c(x)=1\) for \(x\geq c\) and \(F_c(x)=0\) for \(x<c\). It is
      discrete, with \(S=\{c\}\), and it is the only probability measure
      on \(\mathcal B(\mathbb R)\) whose distribution function is
      \(F_c\). For \(c=0\), \(F_0\) is the function of Durrett's
      Example~1.2.8.
\item The following are equivalent: \(\mathbb P_X=\delta_c\);
      \(p_X(c)=1\); \(X=c\) almost surely.
\item Let \(D=\{x\in\mathbb R:p_X(x)>0\}\), an at most countable set.
      Then \(\mathbb P_X\) is discrete if and only if
      \(\sum_{x\in D}p_X(x)=1\), and in that case, for every Borel set
      \(B\),
      \[
       \mathbb P_X(B)=\sum_{x\in D\cap B}p_X(x)
                     =\sum_{x\in D}p_X(x)\,\delta_x(B).
      \]
\end{enumerate}
\end{hblemma}

\section{Source forms of uniform integrability}
\label{sec:p15-w02c-p15-014-ui-source}

\Cref{def:p15-uniform-integrability} defines uniform integrability for
families of integrable variables. Durrett's definition carries no
integrability hypothesis, and his moment criterion assumes only a bound
on \(E\varphi(|X_i|)\). Both are given here in his form, with integrability
of each member derived rather than assumed.

\begin{hbconvention}{def:p15-w02c-p15-014-ui-arbitrary}{Durrett's definition for any collection}
Let \(X_i\), \(i\in I\), be real random variables on
\((\Omega,\mathcal F,P)\), with no integrability assumed. For
\(M\geq0\) the number \(E(|X_i|;|X_i|>M)\) is the integral of a
nonnegative variable and lies in \([0,\infty]\). Durrett calls the
collection \emph{uniformly integrable} if
\[
 \lim_{M\to\infty}\sup_{i\in I}E(|X_i|;|X_i|>M)=0
\]
\cite[sec.~4.6, p.~244]{DurrettPTEAuthor2019}. He then remarks that a
choice of \(M\) making the supremum less than one gives
\(\sup_{i\in I}E|X_i|\leq M+1<\infty\), and that this bound will be used
several times \cite[p.~244]{DurrettPTEAuthor2019}. The remark is what
makes his definition agree with \cref{def:p15-uniform-integrability}:
integrability of every member is a consequence of the displayed
condition, not a separate hypothesis
(\cref{thm:p15-w02c-p15-014-ui-arbitrary}).
\index{uniform integrability!for an arbitrary collection}
\end{hbconvention}

\begin{hblemma}{thm:p15-w02c-p15-014-ui-arbitrary}{proved=annex:p15-w02c-p15-014-ui-arbitrary}{Uniform integrability forces integrability}
\par\nobreak
Let \(X_i\), \(i\in I\), be real random variables, with no integrability
assumed.
\begin{enumerate}
\item If \(M\geq0\) and \(\sup_{i\in I}E(|X_i|;|X_i|>M)\leq1\), then
      \(\sup_{i\in I}E|X_i|\leq M+1\).
\item If some \(X_j\) has \(E|X_j|=\infty\), then
      \(E(|X_j|;|X_j|>M)=\infty\) for every \(M\geq0\), so the collection
      is not uniformly integrable in Durrett's sense.
\item The collection is uniformly integrable in Durrett's sense if and
      only if every \(X_i\) is integrable and the family is uniformly
      integrable in the sense of \cref{def:p15-uniform-integrability}.
\end{enumerate}
\end{hblemma}

\begin{hblemma}{thm:p15-w02c-p15-014-ui-moment-source}{proved=annex:p15-w02c-p15-014-ui-moment}{Durrett's moment criterion with no integrability hypothesis}
\par\nobreak
Let \(\varphi:[0,\infty)\to[0,\infty)\) be a function with
\(\varphi(x)/x\to\infty\) as \(x\to\infty\). Let \(X_i\), \(i\in I\), be
real random variables such that each \(\varphi(|X_i|)\) is a random
variable, and let \(C<\infty\) satisfy \(E\varphi(|X_i|)\leq C\) for all
\(i\in I\). For \(M\geq0\) put
\[
 \varepsilon_M=\sup\{x/\varphi(x):x\geq M\}.
\]
\begin{enumerate}
\item There is \(M_0>0\) with \(\varphi(x)\geq x>0\) for every
      \(x\geq M_0\). For \(M\geq M_0\), \(\varepsilon_M\) is a number in
      \([0,1]\), nonincreasing in \(M\), and \(\varepsilon_M\to0\) as
      \(M\to\infty\).
\item For \(M\geq M_0\) and every \(i\in I\),
      \(E(|X_i|;|X_i|>M)\leq C\varepsilon_M\).
\item Every \(X_i\) is integrable, with \(E|X_i|\leq M_0+C\), and the
      family is uniformly integrable in Durrett's sense and in the sense
      of \cref{def:p15-uniform-integrability}.
\end{enumerate}
\end{hblemma}

This is Durrett's Theorem~4.6.2, which he offers as a common way to check
uniform integrability, with the examples \(\varphi(x)=x^p\) for \(p>1\)
and \(\varphi(x)=x\log^+x\); his proof is the bound of item~(2) with the
same \(\varepsilon_M\) \cite[Theorem~4.6.2, p.~245]{DurrettPTEAuthor2019}.
He asks only that \(\varphi\geq0\) be a function with
\(\varphi(x)/x\to\infty\). His hypothesis \(E\varphi(|X_i|)\leq C\)
presupposes that \(\varphi(|X_i|)\) has an expectation, and the lemma
states that presupposition explicitly; it holds whenever \(\varphi\) is
Borel (\cref{thm:p15-measurable-operations}). The handbook's
\cref{thm:p15-ui-moment-criteria} is the same criterion for Borel
\(\Phi\); item~(3) shows that its conclusion needs no assumption that the
\(X_i\) are integrable, because the moment bound supplies it.
\index{uniform integrability!moment criterion}

\begin{hbexample}{def:p15-w02c-p15-014-ui-moment-example}{The criterion on spikes of height \(\sqrt n\)}
On the uniform space of \cref{thm:p15-lebesgue-unit-interval} let
\(X_n=\sqrt n\,\mathbf1_{(0,1/n)}\) for \(n\geq1\), and take
\(\varphi(x)=x^2\). Then \(E\varphi(X_n)=n\cdot(1/n)=1\), so \(C=1\).
For \(x\geq1\), \(\varphi(x)=x^2\geq x>0\), so \(M_0=1\) serves in
item~(1), and \(\varepsilon_M=\sup\{1/x:x\geq M\}=1/M\). Item~(2) gives
\(E(X_n;X_n>M)\leq1/M\) for \(M\geq1\). Directly, \(E(X_n;X_n>M)\) is
\(\sqrt n/n=n^{-1/2}\) when \(\sqrt n>M\) and zero otherwise, so it is
less than \(1/M\) in every case. Its largest value is taken at the least
integer \(n\) above \(M^2\); that \(n\) is at most \(M^2+1\), so the value
is at least \((M^2+1)^{-1/2}\), and its ratio to the bound \(1/M\) tends
to one as \(M\to\infty\). Item~(3) gives \(E X_n\leq2\);
directly \(E X_n=n^{-1/2}\leq1\).
\end{hbexample}

\begin{hbwarning}{def:p15-w02c-p15-014-ui-moment-limits}{What the moment criterion needs}
\begin{enumerate}
\item The growth condition cannot be weakened to \(\varphi(x)=x\): the
      spikes \(X_n=n\mathbf1_{(0,1/n)}\) have \(E|X_n|=1\) for every
      \(n\) and are not uniformly integrable
      (\cref{thm:p15-ui-small-sets}). For \(\varphi(x)=x\),
      \(\varepsilon_M=1\) for every \(M>0\), and the bound of item~(2) does
      not tend to zero.
\item The bound must be uniform in \(i\). The same spikes have
      \(E\varphi(|X_n|)=n\) for \(\varphi(x)=x^2\), finite for each \(n\)
      but with no common bound \(C\).
\item The criterion is sufficient, not necessary, for a given
      \(\varphi\). On the uniform space, the variable \(X(u)=u^{-1/2}\)
      has \(EX=\int_0^1u^{-1/2}\,du=2\), so the one-member family
      \(\{X\}\) is uniformly integrable (\cref{thm:p15-ui-small-sets}),
      while \(EX^2=\int_0^1u^{-1}\,du=\infty\), so the criterion with
      \(\varphi(x)=x^2\) does not apply to it. The converse in
      \cref{thm:p15-ui-moment-criteria} supplies, for each uniformly
      integrable family, some \(\Phi\) for which the criterion does
      apply.
\end{enumerate}
\end{hbwarning}

\section{Source discussions of conditional limits}
\label{sec:p15-w02c-p15-014-conditional-remarks}

\begin{hbexample}{def:p15-w02c-p15-014-levy-readings}{Obvious or incredible: two readings of L\'evy's zero-one law}
Let \(\mathcal F_n\) increase, \(\mathcal F_\infty=\sigma(\bigcup_n\mathcal F_n)\)
and \(A\in\mathcal F_\infty\). L\'evy's zero-one law,
\cref{thm:p15-levy-information-law}, says
\(\mathbb P(A\mid\mathcal F_n)\to\mathbf1_A\) almost surely. After
stating it, Durrett borrows a line from K.~L.~Chung, who asks the reader
to decide ``whether it is obvious or incredible''
\cite[p.~247]{DurrettPTEAuthor2019}; Durrett gives no reference to a
work of Chung for the line. He then argues for each view.

\emph{The obvious reading.} The indicator \(\mathbf1_A\) is
\(\mathcal F_\infty\)-measurable, and the information \(\mathcal F_n\)
increases to \(\mathcal F_\infty\). The best prediction of
\(\mathbf1_A\) from \(\mathcal F_n\) should therefore approach the best
prediction from \(\mathcal F_\infty\), which is \(\mathbf1_A\) itself
\cite[p.~247]{DurrettPTEAuthor2019}. Each step has a proved counterpart:
\(\mathbb E(\mathbf1_A\mid\mathcal F_\infty)=\mathbf1_A\) because a
measurable integrable variable is its own conditional expectation
(\cref{thm:p15-conditional-basic}), and
\(\mathbb E(\mathbf1_A\mid\mathcal F_n)\to\mathbb E(\mathbf1_A\mid\mathcal F_\infty)\)
almost surely and in \(L^1\) is L\'evy's upward theorem
(\cref{thm:p15-levy-upward}).

\emph{The incredible reading.} Suppose that \(X_1,X_2,\ldots\) are
independent random variables. Put
\(\mathcal F_n=\sigma(X_1,\ldots,X_n)\) and let \(A\) be a tail event.
Then \(A\) is independent of every \(\mathcal F_n\), so
\[
 \mathbb E(\mathbf1_A\mid\mathcal F_n)=\mathbb P(A)\qquad(n\geq1),
\]
a constant. The law says that this constant sequence converges to
\(\mathbf1_A\) almost surely, so \(\mathbb P(A)=\mathbf1_A\) almost
surely, and \(\mathbb P(A)\in\{0,1\}\): Kolmogorov's zero-one law
follows \cite[pp.~247--248]{DurrettPTEAuthor2019}. The complete proof,
including the facts that \(A\in\mathcal F_\infty\) and that \(A\) is
independent of each \(\mathcal F_n\), is the second part of the annex
proof of \cref{thm:p15-levy-information-law}.

\emph{Evaluation.} Durrett concedes that the tail argument may not
establish that the theorem is literally incredible, and judges that it, together with later applications, shows the theorem to
be very useful \cite[p.~248]{DurrettPTEAuthor2019}. One such later
application in this handbook is
\cref{thm:p15-extinction-or-escape}, his Exercises~4.6.4--4.6.5. The
two readings and the evaluation are the source's commentary; the
mathematical content is in the theorems cited.
\end{hbexample}

\begin{hbconvention}{def:p15-w02c-p15-014-dyadic-price}{Step approximation: the martingale route and the direct route}
\Cref{thm:p15-dyadic-martingales} obtains, for integrable \(f\) on
\([0,1)\), the dyadic step functions \(\mathbb E(f\mid\mathcal D_n)\)
converging to \(f\) in \(L^1\), so that every such \(f\) is within any
\(\epsilon>0\) of a step function in \(L^1\). This is Durrett's
Exercise~4.6.3, and he remarks that this approximation is much simpler
than the hands-on argument of his Exercise~1.4.3, although it uses a
great deal of machinery \cite[Exercise~4.6.3, p.~249]{DurrettPTEAuthor2019}.

The direct route of Exercise~1.4.3
\cite[Exercise~1.4.3, p.~24]{DurrettPTEAuthor2019} is proved in this
handbook as part of \cref{thm:p15-criteria-real-integration}: an
integrable function is approximated by a simple function, each set of
the simple function by a finite union of intervals, and the result is a
step function on intervals. The martingale route replaces those
approximations of sets by the existence of conditional expectations
(\cref{thm:p15-conditional-existence}), uniform integrability of a
closed martingale (\cref{thm:p15-ui-conditional-family}) and L\'evy's
upward theorem (\cref{thm:p15-levy-upward}), applied to the dyadic
information of \cref{def:p15-dyadic-information}. Its conclusion is
also more specific: the approximants are the averages of \(f\) over
dyadic intervals, and they converge almost everywhere as well as in
\(L^1\). The comparison of simplicity is Durrett's judgment.
\end{hbconvention}

\begin{hbconvention}{def:p15-w02c-p15-014-conditional-dominated}{The dominated convergence theorem for conditional expectations}
In Durrett's text \cite[Theorem~4.6.10, p.~248]{DurrettPTEAuthor2019}
Theorem~4.6.10 carries the title \emph{Dominated convergence theorem
for conditional expectations}. Its hypotheses are
\(Y_n\to Y\) almost surely, \(|Y_n|\leq Z\) for all \(n\) with
\(EZ<\infty\), and \(\mathcal F_n\uparrow\mathcal F_\infty\); its
conclusion is
\[
 E(Y_n\mid\mathcal F_n)\to E(Y\mid\mathcal F_\infty)\quad\text{a.s.}
\]
In this handbook the theorem is item~(2) of
\cref{thm:p15-changing-conditional-limits}, which adds convergence in
\(L^1\), and its proof is the annex proof of that theorem. The name
records the role of the integrable bound \(Z\), as in the dominated
convergence theorem for expectations
(\cref{thm:p15-dominated-bounded-convergence}). Two features separate
it from that theorem and from the conditional form with fixed
information in \cref{thm:p15-conditional-limits}~(3): the variable and
the information change together, and the limit is taken with respect to
\(\mathcal F_\infty\). With \(\mathcal F_n=\mathcal G\) for every \(n\),
so that \(\mathcal F_\infty=\mathcal G\), it reduces to the almost-sure
part of \cref{thm:p15-conditional-limits}~(3). Without domination the
almost-sure conclusion can fail even for a uniformly integrable
sequence and fixed information; that is Durrett's Example~4.6.11, constructed in
\cref{thm:p15-ui-counterexamples}.
\index{dominated convergence theorem!for conditional expectations}
\end{hbconvention}

\begin{hbexample}{def:p15-w02c-p15-014-prior-posterior}{Prior and posterior in Durrett's Exercise 4.6.1}
\par\nobreak
Durrett's Exercise~4.6.1 takes \(Z_1,Z_2,\ldots\) i.i.d.\ with
\(E|Z_i|<\infty\), a random variable \(\theta\) with finite mean
independent of them, and \(Y_i=Z_i+\theta\), and asks for a proof that
\(E(\theta\mid Y_1,\ldots,Y_n)\to\theta\) almost surely
\cite[Exercise~4.6.1, p.~248]{DurrettPTEAuthor2019}. He gives the
statistical reading of the model when each \(Z_i\) is normal with mean
0 and variance 1: the observations \(Y_1,Y_2,\ldots\) are then drawn
from a normal population whose variance is 1 and whose mean, the value
of \(\theta\), is unknown. In that reading he names two objects
\cite[p.~248]{DurrettPTEAuthor2019}:
\begin{enumerate}
\item the \emph{prior distribution}, which is the law of \(\theta\)
      (\cref{def:p15-pushforward-law});
\item the \emph{posterior distribution} after \(n\) observations, which
      is \(P(\theta\in\cdot\mid Y_1,\ldots,Y_n)\).
\end{enumerate}
The posterior is read here as a regular conditional law of \(\theta\)
given \(\mathcal F_n=\sigma(Y_1,\ldots,Y_n)\)
(\cref{def:p15-regular-conditional-law}), a probability kernel
\(K_n\) with \(K_n(\cdot,A)=P(\theta\in A\mid\mathcal F_n)\) almost surely
for every Borel \(A\). One exists because \(\theta\) is real
(\cref{thm:p15-real-conditional-law}), and any two agree outside one
null set. The \emph{posterior mean} is
\(\int x\,K_n(\cdot,dx)\), the mean of the posterior distribution.

The exercise's conclusion is \cref{thm:p15-posterior-mean-consistency}
in the handbook's notation, with \(\xi_i=Z_i\) and \(\mu=EZ_1\).
\Cref{thm:p15-w02c-p15-014-posterior-mean} restates it for the posterior
distribution: the posterior mean is a version of
\(E(\theta\mid\mathcal F_n)\) and converges to the true value of
\(\theta\). The proof needs no density for the prior and no normality of
the \(Z_i\); normality enters only through the statistical reading. The
statement concerns the posterior mean in this additive model, not the
whole posterior distribution and not other sampling models.
\index{prior distribution}\index{posterior distribution}
\end{hbexample}

\begin{hbcorollary}{thm:p15-w02c-p15-014-posterior-mean}{proved=annex:p15-w02c-p15-014-posterior-mean}{The posterior mean converges to the parameter}
In the setting of \cref{def:p15-w02c-p15-014-prior-posterior}, let
\(K_n\) be a regular conditional law of \(\theta\) given
\(\mathcal F_n=\sigma(Y_1,\ldots,Y_n)\). Then
\(\int|x|\,K_n(\cdot,dx)<\infty\) almost surely, the posterior mean
\(m_n=\int x\,K_n(\cdot,dx)\), set equal to 0 where the absolute integral
is infinite, is a version of \(E(\theta\mid\mathcal F_n)\), and
\[
 m_n\to\theta\quad\text{almost surely and in }L^1 .
\]
\end{hbcorollary}

\begin{hbconvention}{def:p15-w02c-p15-014-reversed-martingales}{Reversed martingales, and why decreasing information makes convergence simple}
Durrett's term is \emph{backwards martingale}: a martingale indexed by
the negative integers, \(X_n\) for \(n\leq0\), adapted to increasing
sigma-fields \(\mathcal F_n\) with \(E(X_{n+1}\mid\mathcal F_n)=X_n\) for
\(n\leq-1\). He notes that some authors call these martingales
\emph{reversed} \cite[sec.~4.7, p.~249]{DurrettPTEAuthor2019}. They are
the backward martingales of \cref{def:p15-backward-martingale}, which
uses nonnegative reverse-time indices \(Y_n=X_{-n}\) and
\(\mathcal G_n=\mathcal F_{-n}\); a backwards, backward or reversed
martingale is the same object.

Durrett justifies giving these processes their own treatment by a claim:
their limit theory is especially simple, and the reason he gives is that
the information shrinks as the index runs down to \(-\infty\)
\cite[sec.~4.7, p.~249]{DurrettPTEAuthor2019}. His proof of the
convergence theorem \cite[Theorem~4.7.1, pp.~249--250]{DurrettPTEAuthor2019}
shows the two places where the decreasing information does the work,
and both appear in the annex proof of
\cref{thm:p15-backward-martingale-convergence}.
\begin{enumerate}
\item \emph{A bound with no hypothesis.} For each \(n\) the finite
      string \(X_{-n},\ldots,X_0\) is a forward martingale ending at the
      fixed variable \(X_0\), so the upcrossing inequality
      (\cref{thm:p15-upcrossing-inequality}) bounds the expected number of
      upcrossings of \([a,b]\) by \(E(X_0-a)^+/(b-a)\), a number that
      does not depend on \(n\). Letting \(n\to\infty\), the total number
      of upcrossings has finite mean, so every backwards martingale has
      an almost-sure limit as \(n\to-\infty\), finite because
      \(E|X_n|\leq E|X_0|\) and Fatou's lemma bound its mean.
\item \emph{A closing variable for free.} The martingale property gives
      \(X_n=E(X_0\mid\mathcal F_n)\) for every \(n\leq0\), so the
      sequence consists of conditional expectations of one integrable
      variable. It is therefore uniformly integrable
      (\cref{thm:p15-ui-conditional-family}), and the convergence holds
      in \(L^1\) (\cref{thm:p15-vitali-convergence}).
\end{enumerate}
For a forward martingale neither step is automatic. Almost-sure
convergence needs a moment bound such as
\(\sup_nE(X_n^+)<\infty\) (\cref{thm:p15-martingale-convergence}), and
even a nonnegative martingale with \(EX_n=1\) for all \(n\) can converge
almost surely to zero without converging in \(L^1\)
(\cref{thm:p15-martingale-source-counterexamples}~(1)). The forward
conclusions of the second step hold only for closed martingales
(\cref{thm:p15-closed-martingale}); in reverse time every martingale is
closed by its last term \(X_0\). That is the content of the claim that
the theory is simple: no condition beyond the definition is needed for
convergence almost surely and in \(L^1\).
\index{martingale!reversed}\index{martingale!backwards}
\end{hbconvention}

\section{Relationships and sources}
\label{sec:p15-w02c-p15-014-relations}

\Cref{thm:p15-w02c-p15-014-probability-mean} refines item~(6) of
\cref{thm:p15-w02b-p15-013-comparison} and uses
\cref{thm:p15-w02b-p15-013-vitali} and the first-mean results
\cref{thm:p15-ui-small-sets,thm:p15-vitali-convergence,thm:p15-probability-l1-tools}.
Its models come from \cref{thm:p15-w02b-p15-013-countermodels}. The
statements of \cref{sec:p15-w02c-p15-014-foundations} add the source's
forms to \cref{def:p15-sigma-algebra}, \cref{thm:p15-measure-laws},
\cref{def:p15-borel-sets},
\cref{def:p15-measurable-map}, \cref{def:p15-pi-lambda},
\cref{thm:p15-measure-uniqueness} and \cref{def:p15-pushforward-law}, and
\cref{thm:p15-w02c-p15-014-borel-functions} extends the continuous case of
\cref{thm:p15-measurable-operations}. The statements of
\cref{sec:p15-w02c-p15-014-ui-source} accompany
\cref{def:p15-uniform-integrability,thm:p15-ui-small-sets,thm:p15-ui-moment-criteria}. Further source forms of the same
foundations are given in Chapter~\ref{sec:p15-w02c-p15-002-source-variants}: Durrett's discrete probability spaces and
the astragalus weights (\cref{def:p15-w02c-p15-002-astragali}); the
open rational rays and the open rectangles with real endpoints as
generators, with the test by \(\{X<q\}\) for rational \(q\)
(\cref{thm:p15-w02c-p15-002-generator-variants}); the notation
\(X\in\mathcal F\) for measurability
(\cref{def:p15-w02c-p15-002-membership-notation}); the indicator
terminology (\cref{def:p15-w02c-p15-002-indicator-terminology});
equality in distribution and its notation
(\cref{def:p15-w02c-p15-002-equal-in-distribution}); the
extended line, random variables on an event and almost sure convergence
without a prescribed limit
(\cref{def:p15-w02c-p15-002-extended-line},
\cref{thm:p15-w02c-p15-002-order-borel},
\cref{thm:p15-w02c-p15-002-domain-extension}); the measurable
convergence set (\cref{thm:p15-w02c-p15-002-convergence-set}); and
uniqueness for measures finite on an exhausting sequence in a pi-system
(\cref{thm:p15-w02c-p15-002-sigma-finite-uniqueness}). The statements of
\cref{sec:p15-w02c-p15-014-conditional-remarks} accompany
\cref{thm:p15-levy-information-law,thm:p15-dyadic-martingales,thm:p15-changing-conditional-limits,thm:p15-posterior-mean-consistency},
and \cref{def:p15-w02c-p15-014-reversed-martingales} accompanies
\cref{def:p15-backward-martingale,thm:p15-backward-martingale-convergence}.

\emph{Sources.} Durrett, \emph{Probability: Theory and Examples},
author version 5 dated 11 January 2019: the definition of a sigma-field
and its stated reason \cite[sec.~1.1, p.~1]{DurrettPTEAuthor2019};
Theorem~1.1.1 and its subadditivity item \cite[p.~2]{DurrettPTEAuthor2019}; the
notation for the Borel sets \cite[p.~3]{DurrettPTEAuthor2019}; the
definition of an algebra or field and Example~1.1.6
\cite[p.~4]{DurrettPTEAuthor2019}; the definition of a random variable
\cite[sec.~1.2, p.~10]{DurrettPTEAuthor2019}; discrete distributions and
Example~1.2.8 \cite[p.~13]{DurrettPTEAuthor2019}; the generality of
measurable maps and random vectors \cite[sec.~1.3, p.~15]{DurrettPTEAuthor2019};
the remark on Theorem~1.3.4 and Theorems~1.3.5--1.3.6
\cite[p.~16]{DurrettPTEAuthor2019};
Exercise~1.4.3 \cite[p.~24]{DurrettPTEAuthor2019}; point masses in
Section~3.2 \cite[p.~117]{DurrettPTEAuthor2019}; the definition of
uniform integrability and the bound that follows it
\cite[sec.~4.6, p.~244]{DurrettPTEAuthor2019}; Theorems~4.6.2
and~4.6.3 \cite[p.~245]{DurrettPTEAuthor2019}; Theorem~4.6.9 and its discussion
\cite[pp.~247--248]{DurrettPTEAuthor2019}; Theorem~4.6.10 and
Exercise~4.6.1 \cite[p.~248]{DurrettPTEAuthor2019}; Exercise~4.6.3 and
its remark \cite[p.~249]{DurrettPTEAuthor2019}; backwards martingales,
the name reversed and the reason for their simple convergence theory
\cite[sec.~4.7, p.~249]{DurrettPTEAuthor2019}; Theorem~4.7.1 and its
proof \cite[pp.~249--250]{DurrettPTEAuthor2019}; and the pi-system
example, the attribution to Dynkin and Theorem~A.1.4
\cite[app.~A.1, p.~456]{DurrettPTEAuthor2019}.
